% mode: ini-etex
% WEB 1435-1438, new_font.  The font identifier's name t: a control word, a character (u-single_base), \csname\endcsname (FONT),
% an active character (FONT and the character, a new string).  "at" with sizes at the bounds, 0, negative, 2048pt; "scaled"
% with 0, 1, 32768, 32769 and negative.  WEB 1438: a font already loaded is found again by name, area and size, in all three
% ways of giving the size, and is not found when the size differs; then the frozen identifier's name is the latest given.  The
% font switch WEB 1395, local and global.  A group undoes the identifier but not the font.  \font\x=\x, keywords after
% the name, an empty name.  No missing fonts here, because that runs mktextfm; see new_font_missing.
\input prelude
\errorcontextlines=0
\def\w#1{\show#1\showthe#1 }
\font\a=cmr10 \w\a \message{\fontname\a}
\font\a=cmr10 at 10pt \w\a
\font\a=cmr10 scaled 1000 \w\a
\font\a=cmr10 at 20pt \w\a
\font\a=cmr10 scaled 2000 \w\a
\font\a=cmr10 scaled 2000 \w\a
\font\a=cmr10 scaled 1095 \w\a
\font\a=cmr10 at 10.95pt \w\a
\font\a=cmr10 AT 12pt \w\a
\font\a=cmr10 SCALED 1200 \w\a
\font\a=cmr10 aT 12pt \w\a
\font\a= cmr10 \w\a
\font\a cmr10 \w\a
\font\a="cmr10" \w\a
\font\a={cmr10} \w\a
\font\a=cmr10.tfm \w\a
\def\n{cmr10}\font\a=\n \w\a
\def\s{at 15pt}\font\a=cmr10 \s \w\a
\font\a=cmr10 \relax \w\a
% at, bounds
\font\a=cmr10 at 2047.99998pt \w\a
\font\a=cmr10 at 2048pt \w\a
\font\a=cmr10 at 2048.5pt \w\a
\font\a=cmr10 at 16383pt \w\a
\font\a=cmr10 at 0pt \w\a
\font\a=cmr10 at -1pt \w\a
\font\a=cmr10 at 1sp \w\a
\font\a=cmr10 at 16384pt \w\a
\font\a=cmr10 at x \w\a x\par
\font\a=cmr10 at \w\a
\dimen1=9pt \font\a=cmr10 at \dimen1 \w\a
\font\a=cmr10 at 2\dimen1 \w\a
% scaled, bounds
\font\a=cmr10 scaled 1 \w\a
\font\a=cmr10 scaled 32768 \w\a
\font\a=cmr10 scaled 32769 \w\a
\font\a=cmr10 scaled 0 \w\a
\font\a=cmr10 scaled -1 \w\a
\font\a=cmr10 scaled 2147483647 \w\a
\font\a=cmr10 scaled x \w\a x\par
\font\a=cmr10 scaled \w\a
\count1=1500 \font\a=cmr10 scaled \count1 \w\a
% partial keywords go back to the input: the text is typeset
\font\a=cmr10 a \w\a \par
\font\a=cmr10 s \w\a \par
\font\a=cmr10 atx \w\a \par
\font\a=cmr10 scal 5 \w\a \par
% names of the identifier
\font\.=cmr10 \w\.
\font\^^M=cmr10 \w\^^M
\font\ =cmr10 \w\ %
\font~=cmr10 \w~
\font\relax=cmr10 \showthe\relax
\let\relax=\undefinedone
\expandafter\font\csname\endcsname=cmr10 \expandafter\w\csname\endcsname
\expandafter\font\csname a b\endcsname=cmr10 \expandafter\w\csname a b\endcsname
\font\new=cmr7 \w\new
\font\newer=cmr7 \w\newer \showthe\new
\showthe\font
\new \showthe\font \show\new
\newer \showthe\font
% an active character that prints as ^^-notation
\catcode`^^A=13 \font^^A=cmr8 \w^^A
\catcode`^^A=15
% local, global
\font\l=cmr9 \begingroup\font\l=cmr12 \w\l\endgroup\w\l
\begingroup\global\font\l=cmr12 \endgroup\w\l
\begingroup\font\m=cmr12 at 13pt \endgroup\show\m
\font\m=cmr12 at 13pt \show\m \showthe\m
\begingroup\font\o=cmr12 at 14pt \global\let\oo=\o \endgroup\show\o \show\oo \showthe\oo
\oo \showthe\font
\globaldefs=1 \begingroup\font\p=cmr12 at 16pt \endgroup\show\p \globaldefs=0
\globaldefs=-1 \begingroup\global\font\q=cmr12 at 17pt \endgroup\show\q \globaldefs=0
% the font selection
\tenrm \showthe\font
\begingroup\new\showthe\font\endgroup\showthe\font
\begingroup\global\new\showthe\font\endgroup\showthe\font
\tenrm
\tracingassigns=1 \tracingrestores=1
\tenrm \tenrm \new \new
\begingroup \tenrm \global\tenrm \endgroup
\font\t=cmr7
\font\t=cmr7
\begingroup\font\t=cmr9\endgroup
\tracingassigns=0 \tracingrestores=0
% the identifier is \nullfont while the name is scanned
\font\z=cmr10 \show\z
\def\getname{\ifx\z\nullfont cmr7\else cmr9\fi}\font\z=\getname \show\z
\font\z=cmr7\z \showthe\font
\tenrm
\nullfont \showthe\font \tenrm
\font\nullfont=cmr10 \showthe\nullfont \show\nullfont
\font\font=cmr10 \show\font
\end
